\chapter{期中综合}

\Needspace{12\baselineskip}

\section{期中综合素养评价（一）}

\Needspace{8\baselineskip}

\pdfbookmark[2]{原书第41页}{orig-41}

\subsection*{原书第41页}

% p041_s01
\BookQuestion{一、1}
要使\(9\times(\square\div4)=72\)，方框里应该填（\hspace{1.8em}）。\par
\begin{BookAnswer}
填\(32\)。先反向求括号里的商：\(72\div9=8\)；再求被除数：\(8\times4=32\)。检验\(9\times(32\div4)=72\)。
\end{BookAnswer}

% p041_s02
\BookQuestion{一、2}
同一平面内，两条直线的位置关系有（\hspace{1.8em}）和（\hspace{1.8em}）。当两条直线相交成直角时，这两条直线互相（\hspace{1.8em}）。\par
\begin{BookAnswer}
依次填“互相平行、相交、垂直”。同一平面内，两条不同直线要么相交，要么平行；相交成直角时，叫互相垂直。
\end{BookAnswer}

% p041_s03
\BookQuestion{一、3（1）}
根据运算律或性质填空：（1）\(357+168+43+32=(357+\underline{\qquad})+(168+\underline{\qquad})\)。\par
\begin{BookAnswer}
依次填\(43,32\)。把\(357\)与\(43\)凑成400，把\(168\)与\(32\)凑成200；调换加数的位置并改变结合顺序，和不变。
\end{BookAnswer}

% p041_s04
\BookQuestion{一、3（2）}
根据运算律或性质填空：（2）\(125\times(80+8)=125\times\underline{\qquad}+125\times\underline{\qquad}\)。\par
\begin{BookAnswer}
依次填\(80,8\)。乘法分配律是\(a\times(b+c)=a\times b+a\times c\)，所以\(125\times(80+8)=125\times80+125\times8\)。
\end{BookAnswer}

% p041_s05
\BookQuestion{一、3（3）}
根据运算律或性质填空：（3）\(a\times99+a=a\times(\underline{\qquad}+\underline{\qquad})\)。\par
\begin{BookAnswer}
依次填\(99,1\)。单独的\(a\)可写成\(a\times1\)，故\(a\times99+a=a\times(99+1)\)。
\end{BookAnswer}

% p041_s06
\BookQuestion{一、4（1）}
（1）一个平行四边形的两条邻边分别长6厘米和4厘米，它的周长是（\hspace{1.8em}）厘米。\par
\begin{BookAnswer}
填\(20\)。相邻边长6厘米、4厘米，另外两边分别与它们相等，周长\((6+4)\times2=20\)厘米。
\end{BookAnswer}

% p041_s07
\BookQuestion{一、4（2）}
（2）一个等边三角形的周长是27厘米，它的边长是（\hspace{1.8em}）厘米。\par
\begin{BookAnswer}
填\(9\)。等边三角形三边相等，边长\(27\div3=9\)厘米。
\end{BookAnswer}

% p041_s08
\BookQuestion{一、5}
在圆圈里填上“＞”“＜”或“＝”。\(125+25\times4\ \bigcirc\ (125+25)\times4\)；\allowbreak{}\(36\times9+36\times2\ \bigcirc\ 36\times9+2\)；\allowbreak{}\(45\times(12+8)\ \bigcirc\ 45\times12+45\times8\)；\allowbreak{}\(600\div(2\times5)\ \bigcirc\ 600\div2\times5\)。\par
\begin{BookAnswer}
依次填\(<,>,=,<\)。\par \(125+25\times4=225\)，\allowbreak{}\((125+25)\times4=600\)，故\(<\)。\par \(36\times9+36\times2=396\)，\allowbreak{}\(36\times9+2=326\)，故\(>\)。\par 乘法分配律给出\(45\times(12+8)=45\times12+45\times8\)，故\(=\)。\par \(600\div(2\times5)=60\)，\allowbreak{}\(600\div2\times5=1500\)，故\(<\)。
\end{BookAnswer}

% p041_s09
\BookQuestion{一、6}
一个直角三角形中，有一个锐角是\(45^\circ\)，另一个锐角是（\hspace{1.8em}），按边分它是（\hspace{1.8em}）三角形。把这样两个完全一样的三角形拼成一个大三角形，大三角形的内角和是（\hspace{1.8em}）；拼成一个正方形，正方形的内角和是（\hspace{1.8em}）。\par
\begin{BookAnswer}
依次填\(45^\circ\)、等腰、\(180^\circ\)、\(360^\circ\)。另一个锐角\(180^\circ-90^\circ-45^\circ=45^\circ\)。两锐角相等，所以是等腰三角形。拼成的大三角形内角和仍为\(180^\circ\)；正方形有4个直角，内角和\(90^\circ\times4=360^\circ\)。
\end{BookAnswer}

\Needspace{8\baselineskip}

\pdfbookmark[2]{原书第42页}{orig-42}

\subsection*{原书第42页}

% p042_s01
\BookQuestion{一、7}
商店购进6箱火龙果，每箱24千克，如果分装成每袋8千克，一共可以分装（\hspace{1.8em}）袋；把这些火龙果按每袋12元售卖，全部售出一共收入（\hspace{1.8em}）元。\par
\begin{BookAnswer}
依次填\(18,216\)。袋数先利用24除以8求每箱袋数，计算更简便；收入按单价乘袋数。\[\begin{aligned}&24\times6\div8\\={}&24\div8\times6\\={}&3\times6\\={}&18\text{(袋)}\\&12\times18\\={}&216\text{(元)}\end{aligned}\]
\end{BookAnswer}

% p042_s02
\BookQuestion{一、8}
右图中，有（\hspace{1.8em}）个三角形，（\hspace{1.8em}）个平行四边形，（\hspace{1.8em}）个梯形。\par
\BookFigure{figures/p042_s02-01.tex}
\begin{BookAnswer}
依次填\(3,3,4\)。三角形有左、中、右3个。平行四边形按两条底边之间的两条同方向斜边配对，共3个。梯形按大小和位置逐一检查，共4个；只数“只有一组对边平行”的图形，不能把平行四边形重复计入。
\end{BookAnswer}

% p042_s03
\BookQuestion{二、1}
与\(720\div(8\times9)\)结果相等的算式是（\hspace{1.8em}）。\par
\noindent A. \(720\div8\times9\)\par
\noindent B. \(720\div8\div9\)\par
\noindent C. \(720\times8\div9\)\par
\noindent D. \(720\div9\times8\)\par
\begin{BookAnswer}
选B。除以两个数的积，等于依次除以这两个数：\(720\div(8\times9)=720\div8\div9\)。
\end{BookAnswer}

% p042_s04
\BookQuestion{二、2}
\(25\times(a+4)\)与\(25\times a+25\)的结果相差（\hspace{1.8em}）。\par
\noindent A. \(24a\)\par
\noindent B. \(a\)\par
\noindent C. 4\par
\noindent D. 75\par
\begin{BookAnswer}
选D。\(25\times(a+4)=25\times a+100\)，另一个式子为\(25\times a+25\)，相差\(100-25=75\)。
\end{BookAnswer}

% p042_s05
\BookQuestion{二、3}
在同一平面内，如果直线\(a\)垂直于直线\(b\)，直线\(b\)垂直于直线\(c\)，那么直线\(a\)和\(c\)（\hspace{1.8em}）。\par
\noindent A. 互相平行\par
\noindent B. 互相垂直\par
\noindent C. 相交但不垂直\par
\noindent D. 无法确定\par
\begin{BookAnswer}
选A。同一平面内，垂直于同一条直线的两条不同直线互相平行，所以\(a\parallel c\)。
\end{BookAnswer}

% p042_s06
\BookQuestion{二、4}
下面每组中的三根小棒，能围成三角形的是（\hspace{1.8em}）（单位：厘米）。\par
\noindent A. 4、5、9\par
\noindent B. 3、4、8\par
\noindent C. 5、6、10\par
\noindent D. 2、2、5\par
\begin{BookAnswer}
选C。检验两短边之和：A为\(4+5=9\)，B为\(3+4<8\)，D为\(2+2<5\)，均不成立；C有\(5+6>10\)，能围成三角形。
\end{BookAnswer}

% p042_s07
\BookQuestion{二、5}
一个三角形的最小角是\(50^\circ\)，这个三角形一定是（\hspace{1.8em}）。\par
\noindent A. 锐角三角形\par
\noindent B. 直角三角形\par
\noindent C. 钝角三角形\par
\noindent D. 无法判断\par
\begin{BookAnswer}
选A。最小角50度，其余角都至少50度，最大角至多\(180^\circ-50^\circ-50^\circ=80^\circ\)，所以一定是锐角三角形。
\end{BookAnswer}

% p042_s08
\BookQuestion{三、1}
直接写得数。\(25\times3\times4=\)，\allowbreak{}\(125\times8=\)，\allowbreak{}\(23\times10=\)，\allowbreak{}\(200-100\div10=\)；\allowbreak{}\(360\div60=\)，\allowbreak{}\(11\times20=\)，\allowbreak{}\((12+8)\times5=\)，\allowbreak{}\(80\div(8\times2)=\)。\par
\begin{BookAnswer}
按题面顺序，得\(300,1000,230,190;6,220,100,5\)。\par \(25\times3\times4=(25\times4)\times3=300\)；\allowbreak{}\(125\times8=1000\)；\allowbreak{}\(23\times10=230\)；\allowbreak{}\(200-100\div10=200-10=190\)。\par \(360\div60=6\)；\allowbreak{}\(11\times20=220\)；\allowbreak{}\((12+8)\times5=20\times5=100\)；\allowbreak{}\(80\div(8\times2)=80\div16=5\)。
\end{BookAnswer}

% p042_s09
\BookQuestion{三、2（1）}
计算下面各题，能简算的要简算。\(640\div[16\times(19-14)]\)。\par
\begin{BookAnswer}
先求最内层的差，再把16与5相乘凑成整十除数80，直接相除。\[\begin{aligned}&640\div[16\times(19-14)]\\={}&640\div[16\times5]\\={}&640\div80\\={}&8\end{aligned}\]
\end{BookAnswer}

% p042_s10
\BookQuestion{三、2（2）}
计算下面各题，能简算的要简算。\(25\times32\times25\)。\par
\begin{BookAnswer}
把32拆成\(4\times8\)，运用乘法结合律，分别凑成整百数。\[\begin{aligned}&25\times32\times25\\&=(25\times4)\times(8\times25)\\&=100\times200\\&=20000.\end{aligned}\]
\end{BookAnswer}

% p042_s11
\BookQuestion{三、2（3）}
计算下面各题，能简算的要简算。\(99\times47+47\)。\par
\begin{BookAnswer}
把最后一个47看作\(1\times47\)，运用乘法分配律，把99个47和1个47合成100个47。\[\begin{aligned}&99\times47+47\\={}&(99+1)\times47\\={}&100\times47\\={}&4700\end{aligned}\]另一写法“先按100个47算，减去1个47后再加回”反映同一凑整结构，不另算一种解法。
\end{BookAnswer}

% p042_s12
\BookQuestion{三、2（4）}
计算下面各题，能简算的要简算。\(35\times98\)。\par
\begin{BookAnswer}
用乘法分配律，把98看作\(100-2\)。\[\begin{aligned}&35\times98\\={}&35\times(100-2)\\={}&35\times100-35\times2\\={}&3500-70\\={}&3430\end{aligned}\]
\end{BookAnswer}

\Needspace{8\baselineskip}

\pdfbookmark[2]{原书第43页}{orig-43}

\subsection*{原书第43页}

% p043_s01
\BookQuestion{四、1}
过直线外的点\(A\)，画已知直线的垂线和平行线。（4分）\par
\BookFigure{figures/p043_s01-01.tex}
\begin{BookAnswer}
先画垂线：把三角尺一条直角边靠在已知直线上，沿直线平移三角尺，使另一条直角边经过\(A\)，再沿这条边画线，垂足处标直角。再借助直尺推移三角尺，过\(A\)画已知直线的平行线。两条红线都要经过\(A\)；平行线与原直线不相交，垂线与原直线交成直角。
\BookFigure{figures/p043_s01-02.tex}
\end{BookAnswer}

% p043_s02
\BookQuestion{四、2}
在下面的方格纸上（每个小方格边长1厘米），画一个底边为5厘米、高为3厘米的平行四边形，并画出它的一条高。（4分）\par
\BookFigure{figures/p043_s02-01.tex}
\begin{BookAnswer}
先在同一格线上取长5格的底边；再在与它相距3格、互相平行的另一格线上取长5格的对边，将对应端点连起来。检查两组对边分别平行，才是平行四边形。用三角尺从上边一点向底边作垂线段，标直角与“高3厘米”。这里3格按垂直方向数，不是沿斜边数。同一底上的高都长3厘米，上边可换一个点另作一条高。
\BookFigure{figures/p043_s02-02.tex}
\end{BookAnswer}

% p043_s03
\BookQuestion{五、1}
工程队修一条公路，原计划每天修120米，8天修完。实际提前2天完成，实际平均每天修多少米？（5分）\par
\begin{BookAnswer}
方法一：路长不变；列式后先算120除以6，凑成整十数20。\par \[\begin{aligned}&120\times8\div(8-2)\\={}&120\times8\div6\\={}&120\div6\times8\\={}&20\times8\\={}&160\text{(米)}\end{aligned}\]答：实际平均每天修160米。\par 方法二：把原来2天的任务平均增加到实际6天，增加部分也先算120除以6。\par \[\begin{aligned}&120+120\times2\div(8-2)\\={}&120+120\times2\div6\\={}&120+120\div6\times2\\={}&120+20\times2\\={}&120+40\\={}&160\text{(米)}\end{aligned}\]答：实际平均每天修160米。
\end{BookAnswer}

% p043_s04
\BookQuestion{五、2}
一个等腰三角形的两条边长分别是6厘米和8厘米。这个等腰三角形的周长可能是多少厘米？（5分）\par
\begin{BookAnswer}
“两条边长6厘米、8厘米”没有指定哪条是底，必须分两种情况。\par 情况一：6厘米为腰，8厘米为底。\[\begin{aligned}&6\times2+8\\&=12+8\\&=20\text{(厘米)}\end{aligned}\]\par 情况二：8厘米为腰，6厘米为底。\[\begin{aligned}&8\times2+6\\&=16+6\\&=22\text{(厘米)}\end{aligned}\]\par 两种都能成三角形，因为\(6+6>8\)、\(8+8>6\)。答：周长可能是20厘米或22厘米。这是两种可能情况，不能只写其中一种。
\end{BookAnswer}

\Needspace{8\baselineskip}

\pdfbookmark[2]{原书第44页}{orig-44}

\subsection*{原书第44页}

% p044_s01
\BookQuestion{五、3}
一根铁丝可以围成一个边长9厘米的正方形，如果改围成一个等边三角形，那么这个等边三角形的边长是多少厘米？（5分）\par
\begin{BookAnswer}
铁丝总长度不变，先按正方形每边9厘米乘4条边列式，再按等边三角形的3条边平均分。计算时先除以3，再乘4，更便于口算。\[\begin{aligned}&9\times4\div3\\={}&9\div3\times4\\={}&3\times4\\={}&12\text{(厘米)}\end{aligned}\]答：等边三角形的边长是12厘米。\par 另一数量理解：每条原边平均分成3小份，共12小份，每条新边取4小份，直接得到\(9\div3\times4\)。它与上面的简算式一致，不重复列第二遍脱式。
\end{BookAnswer}

% p044_s02
\BookQuestion{五、4（1）}
超市举行“酸奶促销活动”：每瓶酸奶6元，买5瓶送1瓶。李阿姨需要买24瓶酸奶送给敬老院。（1）李阿姨实际需要付多少瓶的钱？（5分）\par
\begin{BookAnswer}
“买5送1”表示每组得到6瓶，只付5瓶的钱。\par 方法一：直接求付费瓶数。\[\begin{aligned}&5\times[24\div(5+1)]\\&=5\times[24\div6]\\&=5\times4\\&=20\text{(瓶)}\end{aligned}\]\par 方法二：先求赠送瓶数，再从总瓶数中扣去。\[\begin{aligned}&24-24\div6\\&=24-4\\&=20\text{(瓶)}\end{aligned}\]\par 答：实际要付20瓶的钱。
\end{BookAnswer}

% p044_s03
\BookQuestion{五、4（2）}
超市举行“酸奶促销活动”：每瓶酸奶6元，买5瓶送1瓶。李阿姨需要买24瓶酸奶送给敬老院。（2）她一共要付多少元？（5分）\par
\begin{BookAnswer}
方法一：每瓶价格乘付费瓶数。\[\begin{aligned}&6\times20\\={}&120\text{(元)}\end{aligned}\]答：一共要付120元。\par 方法二：从24瓶全价扣去赠品价格，反用乘法分配律先算付费瓶数。\[\begin{aligned}&6\times24-6\times(24\div6)\\={}&6\times24-6\times4\\={}&6\times(24-4)\\={}&6\times20\\={}&120\text{(元)}\end{aligned}\]答：一共要付120元。
\end{BookAnswer}

% p044_s04
\BookQuestion{五、5}
实验小学开展“废纸回收，保护环境”活动。四年级回收废纸275千克，五年级回收的废纸比四年级的2倍少50千克。每回收25千克废纸可以挽救1棵大树。两个年级一共可以挽救多少棵大树？（6分）\par
\begin{BookAnswer}
方法一：先求五年级质量，按总回收量列式，再把275与500分别除以25。\par \[\begin{aligned}&[275+(275\times2-50)]\div25\\={}&[275+(550-50)]\div25\\={}&[275+500]\div25\\={}&275\div25+500\div25\\={}&11+20\\={}&31\text{(棵)}\end{aligned}\]答：两个年级一共可以挽救31棵大树。\par 方法二：先按四年级的3份算，再扣去少回收50千克对应的棵数。\par \[\begin{aligned}&275\div25\times3-50\div25\\={}&11\times3-2\\={}&33-2\\={}&31\text{(棵)}\end{aligned}\]答：两个年级一共可以挽救31棵大树。\par 方法三：两个年级分别换算棵数，再相加。\par \[\begin{aligned}&275\div25+(275\times2-50)\div25\\={}&11+(550-50)\div25\\={}&11+500\div25\\={}&11+20\\={}&31\text{(棵)}\end{aligned}\]答：两个年级一共可以挽救31棵大树。
\end{BookAnswer}

\Needspace{12\baselineskip}

\section{期中综合素养评价（二）}

\Needspace{8\baselineskip}

\pdfbookmark[2]{原书第45页}{orig-45}

\subsection*{原书第45页}

% p045_s01
\BookQuestion{一、1}
给算式\(200+370-295\div5\)添上括号，使它的运算顺序变成先算减法，再算加法，最后算除法：（\hspace{1.8em}）。\par
\begin{BookAnswer}
填\([200+(370-295)]\div5\)。题目要求先算减法，把\(370-295\)括起；再与200相加，把这个和整体放进中括号；最后除以5。
\end{BookAnswer}

% p045_s02
\BookQuestion{一、2}
装修公司要给一个面积为96平方米的会议厅铺防滑地砖，每平方米的地砖和施工费用合计48元，完成这个会议厅的地砖铺设总共需要花费（\hspace{1.8em}）元。\par
\begin{BookAnswer}
填\(4608\)。每平方米48元，96平方米的费用为\(48\times96=48\times(100-4)=4800-192=4608\)元。乘法顺序按“每平方米费用×平方米数”。
\end{BookAnswer}

% p045_s03
\BookQuestion{一、3}
钟面上（\hspace{1.8em}）时整和（\hspace{1.8em}）时整，时针和分针互相垂直。\par
\begin{BookAnswer}
依次填\(3,9\)。整时的分针指向12。时针指向3或9时，与12相差3个大格，即\(30^\circ\times3=90^\circ\)，所以互相垂直。
\end{BookAnswer}

% p045_s04
\BookQuestion{一、4}
小新用加法运算律写了一个等式：\(\square2+4\square=\square8+5\square\)。如果这两个加数的和是100，那么这两个加数分别是（\hspace{1.8em}）和（\hspace{1.8em}）。\par
\begin{BookAnswer}
两个加数为\(52,48\)，四个方框依次填\(5,8,4,2\)。算式应用加法交换律，故\(\square2\)与\(5\square\)相同、\(4\square\)与\(\square8\)相同，先得到52和48。检验\(52+48=100\)。
\end{BookAnswer}

% p045_s05
\BookQuestion{一、5（1）}
根据运算律填一填。（1）\(175+264+225+136=(\square\bigcirc\square)\bigcirc(\square\bigcirc\square)\)。\par
\begin{BookAnswer}
一种填写为\((175+225)+(264+136)\)。把175和225凑成400，把264和136也凑成400，应用加法交换律和结合律，总和为800。两组的前后互换也正确。
\end{BookAnswer}

% p045_s06
\BookQuestion{一、5（2）}
根据运算律填一填。（2）\(\square\times99+\square\times\square=a\times(\square\bigcirc1)\)。\par
\begin{BookAnswer}
依次填\(a,a,1,99,+\)。式子为\(a\times99+a\times1=a\times(99+1)\)，这是乘法分配律的逆用。
\end{BookAnswer}

% p045_s07
\BookQuestion{一、6}
印刷厂工人每小时能印刷300份报纸，每天工作8小时，一周（5个工作日）能印刷（\hspace{1.8em}）份报纸。\par
\begin{BookAnswer}
填\(12000\)。每小时300份，按“每份数\(\times\)份数”列式，用乘法结合律先求总工作时间：\[\begin{aligned}&300\times8\times5\\={}&300\times(8\times5)\\={}&300\times40\\={}&12000\text{(份)}\end{aligned}\]一周只算5个工作日。
\end{BookAnswer}

% p045_s08
\BookQuestion{一、7}
在圆圈里填上“＞”“＜”或“＝”。\(38\times(5+6)\ \bigcirc\ 38\times5+6\)，\allowbreak{}\(420\div21\ \bigcirc\ 420\div7\div3\)；\allowbreak{}\(78\times199\ \bigcirc\ 78\times200-1\)，\allowbreak{}\(74\times10\times2\ \bigcirc\ 74\times12\)；\allowbreak{}\(65\times24\ \bigcirc\ 65\times4\times20\)，\allowbreak{}\(107\times5\ \bigcirc\ 105\times7\)。\par
\begin{BookAnswer}
依次填\(>,=,<,>,<,<\)。\par \(38\times(5+6)=418>196=38\times5+6\)。\par \(420\div21=420\div7\div3=20\)。\par \(78\times199=15522<15599=78\times200-1\)，少的是78而不是1。\par \(74\times10\times2=1480>888=74\times12\)。\par \(65\times24=1560<65\times4\times20=5200\)。\par \(107\times5=535<735=105\times7\)。
\end{BookAnswer}

% p045_s09
\BookQuestion{一、8}
斑马线中的每一组白线可以看作互相（\hspace{1.8em}）的线段。\par
\begin{BookAnswer}
填“平行”。一组斑马线中的白线沿同一方向、不会相交，可看作互相平行的线段。
\end{BookAnswer}

% p045_s10
\BookQuestion{一、9}
一个等腰三角形的木板，顶角是底角的4倍，底角是（\hspace{1.8em}）°。\par
\begin{BookAnswer}
填\(30^\circ\)。把一个底角看成1份，另一个底角也是1份，顶角为4份，共6份。每份\(180^\circ\div(1+1+4)=30^\circ\)。
\end{BookAnswer}

\Needspace{8\baselineskip}

\pdfbookmark[2]{原书第46页}{orig-46}

\subsection*{原书第46页}

% p046_s01
\BookQuestion{一、10}
在直角三角形\(ABC\)中（单位：厘米），\allowbreak{}\(AC\)边上的高是（\hspace{1.8em}）厘米，\allowbreak{}\(AB\)边上的高是（\hspace{1.8em}）厘米。三角形（\hspace{1.8em}）中（\hspace{1.8em}）边上的高是1.9厘米。\par
\BookFigure{figures/p046_s01-01.tex}
\begin{BookAnswer}
依次填\(1.8,4,BCD,BC\)。先找底，再找“从对面顶点到这条底所在直线”的垂线段，不能看哪一条边画得较斜或较短。\par 以\(AC\)为底，高是\(BD\)，按标注为1.8厘米；以\(AB\)为底，高是与它垂直的\(BC\)，长4厘米。直角三角形的两条直角边可以互为底和高。\(D\)到\(BC\)的垂线段标为1.9厘米，因此在三角形\(BCD\)中，以\(BC\)为底，对应高按图为1.9厘米。\par 原图的尺寸不能同时成立：\(D\)到\(BC\)的垂直线段长1.9厘米，却比从\(D\)到直线上一点\(B\)的线段\(DB=1.8\)厘米还长，违反“垂线段最短”。以上只按题图标示识别底高关系，不能据这些冲突尺寸另作精确图。
\end{BookAnswer}

% p046_s02
\BookQuestion{二、1}
下列算式中正确的是（\hspace{1.8em}）。\par
\noindent A. \(480\div12\times4=480\div(12\times4)\)\par
\noindent B. \(630\div(7\times9)=630\div7\times9\)\par
\noindent C. \(25\times(40+8)=25\times40+8\)\par
\noindent D. \(900\div25\div4=900\div(25\times4)\)\par
\begin{BookAnswer}
选D。\(900\div25\div4=900\div(25\times4)=9\)，符合连续除法的性质。A把“除后乘4”改成“再除4”；B同样把除以9改成乘9；C漏乘了括号里的8。
\end{BookAnswer}

% p046_s03
\BookQuestion{二、2}
一幢大楼有25层，每层有18扇窗户，每扇窗户有4块玻璃，这幢大楼一共有多少块玻璃？下面列式正确的是（\hspace{1.8em}）。\par
\noindent A. \(25\times(18+4)\)\par
\noindent B. \(25\times18\times4\)\par
\noindent C. \(25\times18-4\)\par
\noindent D. \(25\times18\div4\)\par
\begin{BookAnswer}
选B。每扇4块，先求18扇一层的玻璃数，再求25层：\[\begin{aligned}&4\times18\times25\\={}&(4\times25)\times18\\={}&100\times18\\={}&1800\text{(块)}\end{aligned}\]选项\(25\times18\times4\)与此式由乘法交换律相等。
\end{BookAnswer}

% p046_s04
\BookQuestion{二、3}
下列各组小棒中（单位：厘米），能围成三角形的是（\hspace{1.8em}）。\par
\noindent A. 3、4、7\par
\noindent B. 2、3、6\par
\noindent C. 6、6、10\par
\noindent D. 5、2、2\par
\begin{BookAnswer}
选C。C中两短边之和\(6+6=12>10\)，能围成三角形。A为\(3+4=7\)，B为\(2+3<6\)，D为\(2+2<5\)，均不行。
\end{BookAnswer}

% p046_s05
\BookQuestion{二、4}
\(34\times(\triangle+\square)=34\times\triangle+34\times\square\)，这是运用了（\hspace{1.8em}）。\par
\noindent A. 加法交换律\par
\noindent B. 乘法结合律\par
\noindent C. 乘法分配律\par
\noindent D. 乘法交换律\par
\begin{BookAnswer}
选C。\(34\times(\triangle+\square)=34\times\triangle+34\times\square\)，把34分别乘两个加数，应用的是乘法分配律。
\end{BookAnswer}

% p046_s06
\BookQuestion{二、5}
在梯形里画一条线段，把梯形分成两个图形，这两个图形不可能是（\hspace{1.8em}）。\par
\noindent A. 两个梯形\par
\noindent B. 两个三角形\par
\noindent C. 两个平行四边形\par
\noindent D. 一个三角形和一个平行四边形\par
\begin{BookAnswer}
选C。一刀分成两个平行四边形，会使外部图形也具有两组对边分别平行，与原梯形“只有一组对边平行”矛盾。因此不可能分成两个平行四边形；其余三类均能画出实例。
\end{BookAnswer}

% p046_s07
\BookQuestion{三、1}
直接写得数。\(12\times9=\)，\allowbreak{}\(152\times4=\)，\allowbreak{}\(600\div25=\)，\allowbreak{}\((12+8)\times5=\)，\allowbreak{}\(397-23\times6=\)；\allowbreak{}\(125\times8=\)，\allowbreak{}\(24\times5=\)，\allowbreak{}\(480\div80\div6=\)，\allowbreak{}\(49\times31\approx\)，\allowbreak{}\(61\times95\approx\)。\par
\begin{BookAnswer}
第一行依次为\(108,608,24,100,259\)。\par 第二行依次为\(1000,120,1,1500,6000\)。\par \(12\times9=108\)，\allowbreak{}\(152\times4=608\)，\allowbreak{}\(600\div25=600\times4\div100=24\)，\allowbreak{}\((12+8)\times5=100\)。\par \(397-23\times6=397-138=259\)；\allowbreak{}\(125\times8=1000\)，\allowbreak{}\(24\times5=120\)，\allowbreak{}\(480\div80\div6=6\div6=1\)。\par 估算：\(49\times31\approx50\times30=1500\)，\allowbreak{}\(61\times95\approx60\times100=6000\)。估算写约等号，不能把估计值当精确值。
\end{BookAnswer}

% p046_s08
\BookQuestion{三、2（1）}
计算下面各题，能简算的要简算。\(480-12\times(360\div18)\)。\par
\begin{BookAnswer}
先算括号内除法。\[\begin{aligned}&480-12\times(360\div18)\\&=480-12\times20\\&=480-240\\&=240.\end{aligned}\]
\end{BookAnswer}

% p046_s09
\BookQuestion{三、2（2）}
计算下面各题，能简算的要简算。\(45\times[(63+177)\div24]\)。\par
\begin{BookAnswer}
按括号层次计算。\[\begin{aligned}&45\times[(63+177)\div24]\\&=45\times[240\div24]\\&=45\times10\\&=450.\end{aligned}\]
\end{BookAnswer}

% p046_s10
\BookQuestion{三、2（3）}
计算下面各题，能简算的要简算。\(136\times101-136\)。\par
\begin{BookAnswer}
把单独的136看作\(136\times1\)，运用乘法分配律逆用凑出100。\[\begin{aligned}&136\times101-136\\={}&136\times(101-1)\\={}&136\times100\\={}&13600\end{aligned}\]展开成“13600加136再减136”是同一凑整结构的另写法，不另列一种解法。
\end{BookAnswer}

% p046_s11
\BookQuestion{三、2（4）}
计算下面各题，能简算的要简算。\(125\times48\)。\par
\begin{BookAnswer}
方法一：拆48为\(8\times6\)，运用乘法结合律。\[\begin{aligned}&125\times48\\={}&125\times(8\times6)\\={}&(125\times8)\times6\\={}&1000\times6\\={}&6000\end{aligned}\]\par 方法二：把48看作\(50-2\)，用乘法分配律。\[\begin{aligned}&125\times48\\={}&125\times(50-2)\\={}&125\times50-125\times2\\={}&6250-250\\={}&6000\end{aligned}\]
\end{BookAnswer}

\Needspace{8\baselineskip}

\pdfbookmark[2]{原书第47页}{orig-47}

\subsection*{原书第47页}

% p047_s01
\BookQuestion{四、1}
小区要在一条笔直的公路旁建一个快递驿站，驿站的位置在点\(A\)处（如下图）。现在要从驿站修一条最短的小路连接到公路上，方便居民取快递，请你画出这条小路。（4分）\par
\BookFigure{figures/p047_s01-01.tex}
\begin{BookAnswer}
把三角尺一条直角边靠在公路近侧边缘所在直线上，平移到另一条直角边经过\(A\)，从\(A\)沿这条边画到公路边缘，垂足处标直角。点到直线的垂线段最短，所以这是最短小路。小路的终点是公路边缘，不必延伸到公路中心或另一侧。
\BookFigure{figures/p047_s01-02.tex}
\end{BookAnswer}

% p047_s02
\BookQuestion{四、2}
用直尺和圆规画一个以线段\(AB\)为边的等边三角形。（4分）\par
\BookFigure{figures/p047_s02-01.tex}
\begin{BookAnswer}
圆规开口取\(AB\)的长度，针尖先固定在\(A\)画弧，再保持同样开口、固定在\(B\)画弧，两弧交于\(C\)。连接\(AC,BC\)，并保留交点附近的两段弧。因为\(AC\)与\(BC\)都是按同一个长度\(AB\)画出的，所以\(AB=AC=BC\)，得到等边三角形。若两弧暂未相交，应延长弧线，不能任意改变圆规开口。交点取在\(AB\)的另一侧同样成立。
\BookFigure{figures/p047_s02-02.tex}
\end{BookAnswer}

% p047_s03
\BookQuestion{五、1}
求出\(\angle1\)的度数。（5分）\par
\BookFigure{figures/p047_s03-01.tex}
\begin{BookAnswer}
方法一：先用平角求\(\angle2\)，再用三角形内角和求\(\angle1\)。\[\begin{aligned}\angle2&=180^\circ-138^\circ\\&=42^\circ,\\\angle1&=180^\circ-61^\circ-42^\circ\\&=119^\circ-42^\circ\\&=77^\circ.\end{aligned}\]\par 方法二：把三角形的三个内角与底边上的平角比一比，去掉两边共同的\(\angle2\)，可知\(\angle1\)和\(61^\circ\)合起来等于\(138^\circ\)。\[\begin{aligned}\angle1&=138^\circ-61^\circ\\&=77^\circ.\end{aligned}\]\par 答：\(\angle1=77^\circ\)。第二法的依据仍是已学的“内角和180度”与“平角180度”。
\end{BookAnswer}

% p047_s04
\BookQuestion{五、2}
两架客机从同一机场向相反方向起飞，甲机巡航速度为800千米/时，乙机巡航速度为850千米/时。两架飞机同时飞出2小时后，它们之间相距多少千米？（先画图整理条件和问题，再解答）（6分）\par
\begin{BookAnswer}
画图：以机场为共同起点，向相反方向画两段路程；甲段标\(800\times2\)，乙段标\(850\times2\)，两段合起来就是相距的路程。\par 方法一：先求速度和，再乘共同时间。\[\begin{aligned}&(800+850)\times2\\&=1650\times2\\&=3300\text{(千米)}\end{aligned}\]\par 方法二：分别求两架飞机的路程，再相加。\[\begin{aligned}&800\times2+850\times2\\&=1600+1700\\&=3300\text{(千米)}\end{aligned}\]\par 答：2小时后两架飞机之间相距3300千米。
\BookFigure{figures/p047_s04-01.tex}
\end{BookAnswer}

\Needspace{8\baselineskip}

\pdfbookmark[2]{原书第48页}{orig-48}

\subsection*{原书第48页}

% p048_s01
\BookQuestion{五、3}
猕猴桃每100克大约含62毫克维生素C。成年人每日需摄入超过100毫克维生素C。李叔叔今天吃了200克猕猴桃，能满足一天的维生素C摄入量吗？（6分）\par
\begin{BookAnswer}
每100克大约含62毫克，200克里有\(200\div100=2\)份。\[\begin{aligned}&62\times(200\div100)\\={}&62\times2\\={}&124\text{(毫克)}\end{aligned}\]也可以用重复加法表示两份：\(62+62=124\text{(毫克)}\)，这是相同数量关系的两种表达，不另算独立思路。按题目给定的含量估计，约摄入124毫克维生素C，\allowbreak{}\(124>100\)。答：能满足题目所给的一天维生素C摄入量要求。
\end{BookAnswer}

% p048_s02
\BookQuestion{五、4（1）}
学校有三个热门社团，航模社团报名128人，编程社团报名275人，书法社团报名172人。（1）三个社团一共有多少人报名？（4分）\par
\begin{BookAnswer}
把航模、编程、书法三个社团人数相加，用加法交换律和结合律先凑300：\[\begin{aligned}&128+275+172\\={}&(128+172)+275\\={}&300+275\\={}&575\text{(人)}\end{aligned}\]答：三个社团一共有575人报名。
\end{BookAnswer}

% p048_s03
\BookQuestion{五、4（2）}
学校有三个热门社团，航模社团报名128人，编程社团报名275人，书法社团报名172人。（2）编程社团的报名人数，比航模和书法两个社团报名人数的总和少多少人？（4分）\par
\begin{BookAnswer}
方法一：先求航模和书法两社团总人数，再与编程社团比较。\[\begin{aligned}&128+172-275\\&=300-275\\&=25\text{(人)}\end{aligned}\]\par 方法二：先用编程人数减去书法人数，余下部分再与航模人数比较。\[\begin{aligned}&128-(275-172)\\&=128-103\\&=25\text{(人)}\end{aligned}\]\par 答：编程社团比另外两个社团人数的总和少25人。
\end{BookAnswer}

% p048_s04
\BookQuestion{五、5}
国庆节期间，学校组织四年级90名学生和20位老师去公园玩。用3000元购票后还剩多少元？（6分）\par
\par\smallskip\Needspace{4\baselineskip}\begingroup\small\setlength{\tabcolsep}{3pt}
\noindent\textbf{公园门票收费标准}\par\smallskip
\noindent\begin{tabularx}{\linewidth}{|*{3}{>{\centering\arraybackslash}X|}}\hline
季节 & 成人 & 儿童\\\hline
旺季 & 30元 & 20元\\\hline
淡季 & 15元 & 10元\\\hline
\end{tabularx}\par
\noindent 旺季：4月1日～10月31日\par
\noindent 淡季：11月1日～3月31日\par
\endgroup\smallskip
\begin{BookAnswer}
国庆节在10月，按旺季票价：儿童每人20元，成人每人30元。\par 方法一：先求两类票总费用，再从3000元中扣去。\[\begin{aligned}&3000-(20\times90+30\times20)\\={}&3000-(1800+600)\\={}&3000-2400\\={}&600\text{(元)}\end{aligned}\]答：还剩600元。\par 方法二：先买儿童票，再买成人票，依次扣款。\[\begin{aligned}&3000-20\times90-30\times20\\={}&3000-1800-600\\={}&1200-600\\={}&600\text{(元)}\end{aligned}\]答：还剩600元。\par 方法三：先按110人全买儿童票，再补老师的差价。\[\begin{aligned}&3000-20\times(90+20)-(30-20)\times20\\={}&3000-20\times110-10\times20\\={}&3000-2200-200\\={}&800-200\\={}&600\text{(元)}\end{aligned}\]答：还剩600元。
\end{BookAnswer}
